% =====================================================================
%  Sezione 2.7 — L'altezza del bacino dal pavimento
%  Segue la sezione sulle grandezze integrali.
% =====================================================================

\section{L'altezza del bacino dal pavimento}
\label{sec:pavimento}

La rappresentazione riferita al corpo descritta nella
Sezione~\ref{sec:rappresentazioni} normalizza ogni grandezza rispetto al
bacino e alla lunghezza del tronco. È il suo pregio — rende le grandezze
indipendenti dalla distanza dalla telecamera e dalla corporatura — ma comporta
una conseguenza che è opportuno rendere esplicita: \textbf{cancella per
costruzione l'informazione su dove il corpo si trovi nella stanza}.

\subsection{Perché serve}

Il modello, disponendo delle sole grandezze riferite al corpo, apprende una
regola posturale: un tronco orizzontale implica una caduta. È una regola
corretta su un dataset in cui nessuno si sdraia su un divano, e sbagliata in
un'abitazione reale, dove la posizione orizzontale è una condizione ordinaria.

La misura rende la differenza quantitativa. Fra i fotogrammi negativi in cui il
corpo è orizzontale, l'altezza mediana del bacino vale $+0{,}02$ lunghezze di
torso nel dataset di laboratorio e $+1{,}47$ nelle registrazioni domestiche.
Nel primo caso la persona è \emph{a terra} — sdraiarsi sul pavimento è una delle
attività registrate come normali — nel secondo è su un supporto rialzato. La
grandezza «corpo orizzontale» non distingue nulla nel primo insieme e distingue
tutto nel secondo.

\subsection{Come viene stimata}

Il piano del pavimento non viene fornito da nessuno: il sistema lo stima da sé,
osservando che quando una persona è in piedi il bordo inferiore della sua sagoma
coincide con il pavimento a quella distanza dalla telecamera.

Si selezionano quindi i fotogrammi in cui il soggetto risulta eretto — rapporto
altezza/larghezza del riquadro superiore a una soglia, angolo del tronco
prossimo alla verticale, posa valida — e si raccolgono le coppie (posizione
orizzontale del bacino, bordo inferiore del riquadro). Su queste si adatta una
retta

\[
y_{\text{pav}}(x) = a\,x + c,
\]

che rappresenta il pavimento nell'immagine, prospettiva compresa: in una ripresa
non frontale il pavimento non è una riga orizzontale, e la pendenza $a$ ne tiene
conto.

L'altezza del bacino è quindi la distanza verticale dal piano stimato,
normalizzata sulla lunghezza del tronco:

\[
h_{\text{pav}} = \frac{y_{\text{pav}}(x_{\text{bacino}}) - y_{\text{bacino}}}
                      {L_{\text{tronco}}} .
\]

La normalizzazione rende la grandezza adimensionale e quindi confrontabile fra
stanze diverse e fra persone di corporatura diversa. Ciò che resta specifico
dell'installazione è soltanto il riferimento, cioè i coefficienti $a$ e $c$, e
il sistema li stima senza chiedere nulla all'utente.

\subsection{Un limite dipendente dal punto di vista}

La grandezza funziona però soltanto se la telecamera vede il pavimento come un
piano distinto dalla posizione della persona, cioè se è collocata
approssimativamente all'altezza del busto.

La misura sulle due telecamere del dataset lo mostra con chiarezza. Sulla vista
laterale la separazione fra le mediane di caduta e non-caduta vale
$+1{,}43$ lunghezze di torso; sulla vista frontale rialzata vale $-0{,}18$, cioè
è assente e di segno opposto. Una telecamera zenitale non distingue chi è a
terra da chi è in piedi, perché entrambi si proiettano nella stessa regione
dell'immagine.

L'altezza dal pavimento è quindi una grandezza \emph{condizionata
all'installazione}, e va dichiarata come tale: aggiunge informazione quando la
telecamera è montata di profilo, non ne aggiunge — e può fuorviare — quando è
montata in alto.

\subsection{Una conferma arrivata per altra via}

Un riscontro indipendente di quanto questa informazione conti è emerso dalla
verifica sulle occlusioni riportata nel Capitolo~\ref{cap:sperimentale}.
Mascherando i \emph{landmark} delle anche — cioè proprio il riferimento rispetto
a cui la rappresentazione riferita al corpo normalizza tutto — il risultato non
peggiora ma \textbf{migliora} lievemente.

La spiegazione è che, con il bacino congelato in una posizione fissa, l'altezza
della testa smette di essere misurata rispetto al bacino e viene misurata
rispetto a un punto fisso nella stanza: la rappresentazione riferita al corpo si
trasforma, di fatto, in una rappresentazione riferita all'ambiente. Con
telecamera fissa, un riferimento fisso è più informativo di uno che si sposta
insieme alla persona.

È la stessa informazione che $h_{\text{pav}}$ fornisce esplicitamente, ottenuta
per caso da un esperimento che aveva un altro scopo.
